\chapter{Lexical Derivation and Nominalisation}\label{ch:derivational}\label{ch:derivation-nominalisation}

Double brackets identify lexical slots in the schematic examples
(§\,\ref{sec:lexical-slot-schemata}); a subscript gives the complete
constructional and inflectional cell, including any derivation named there.
Thus \textsc{lex.agt} identifies a lexical agent noun, whereas
\Antip{}.\Pf{}.\Nmlz{} identifies a perfective participant nominal
built on antipassive Voice. The bracketed meaning specifies no stem class
or surface allomorph.

\section{Nominal derivation}

\subsection{Diminutives and augmentatives}

Quantity can be part of a derivational stem alternation. The diminutive of
\ird{stol} is \ird{stólek}, with strong \ird{stól-}; its bound stem is
\ird{stólk-}. Selection of G1 and loss of the fleeting \ird{e} are separate
operations, so bound-stem reduction preserves the long root vowel.

\subsection{Agent nouns: \ird{-(h)ár} and \ird{-ist}}\label{sec:lexical-agent-nouns}

Agent nouns in \ird{-(h)ár} and \ird{-ist} denote the habitual or
professional performer of an action; both suffixes are productive.

Derivational quantity need not occupy the same position in every construction.
Schematic, unattested \ird{mel-} illustrates root-strong iterative
\ird{mél-} versus suffix-strong agent \ird{melár}. These are constructed
illustrations of the quantity contrast, not lexical attestations or evidence
for an unrestricted iterative class.
The comparison concerns construction-specific quantity: vowel alternation
need not follow automatically from the segmental shape of a formant
\cite{akelaitiene1995}. This comparative observation does not establish the
Iridian examples as attested lexemes.

Participant nominals in \ird{-ou}, by comparison, characterise an individual
through an aspectually construed event (§\,\ref{sec:nominalised}). With
transitive bases, agent-oriented forms require antipassive Voice in the
perfective and retrospective, and active Voice in the other aspects.


A characteristic occupation and participation in a particular event can be
predicated of the same individual without making the two descriptions
coextensive. Lexical \ird{marár} names the occupation, whereas
\ird{maržikou} identifies the actor of a bounded cooking event:

\pex\label{ex:deriv-slot-occupation}
\a
\begingl
\gla Pres marár.//
\glb man.\Dir{} cook-\textsc{lex.agt}.\Dir{}//
\glft \trsl{The man is a cook.}
\lingcontext{the man's occupation; no particular cooking occasion is identified}//
\endgl
\a
\begingl
\gla Pres maržikou.//
\glb man.\Dir{} cook-\Antip{}-\Pf{}-\Nmlz{}.\Dir{}//
\glft \trsl{The man is the one who cooked.}
\lingcontext{a completed cooking occasion identifies the man; no occupation is attributed}//
\endgl
\xe

The first predicate classifies the man by a characteristic activity,
whereas the second identifies a participant in the contextually selected
event. A professional cook need not have prepared the meal in question,
while the person who did so need not be a cook by occupation. The
antipassive-based participant predicate licenses no overt patient in its
verbal domain; its nominal dependents have the separate licensing described
in §\,\ref{sec:nm-derived-gerund}.

\subsection{Abstract nouns: \ird{-(i)žnám}}

Abstract formation has two lexical quantity classes. The root-strong class
selects G1 in its licensed base domain and the short suffix
\ird{-(i)žnam}; the suffix-strong class selects the G0 base and long
\ird{-(i)žnám}. Schematically, their structures are
\ird{X}\textsubscript{G1}\ird{-(i)žnam} and
\ird{X}\textsubscript{G0}\ird{-(i)žnám}, respectively. These schematic
forms illustrate class architecture, not attested words. Each derivative's
principal parts record its class and selected stem; the suffix's segmental
shape alone does not assign an unfamiliar base to either class.

The selected derivational grade persists into subsequent inflection unless
an independently licensed weak selector changes it in its own domain. This
principle applies to derivational grades generally: inflection does not
automatically restore the underived base. The short and long abstract
suffixes are construction-specific allomorphs, not instances of unrestricted
rhythmic shortening. Outer anchoring retains its independently specified
quantity, including after a long root or abstract suffix.


\section{Verbal derivation}

Unlike the valency operations described in Chapter~\ref{ch:argument-structure}, the
processes treated here do not occupy obligatory positions in the finite verbal
template. They form specialised lexical stems or stem variants, which then
take the ordinary inflectional morphology of the language. In the modern
language, denominal verbalisation and a restricted class of lexical preverbs
form new verbal stems.

Lexically specified intensive derivation can likewise select G1: \ird{muž-}
has intensive \ird{múž-}. The attributive quantity classes distinguish
root-strong \ird{múžna} from suffix-strong \ird{hradná}. The long vowel
marks morphological strength rather than stress, which remains initial in
the lexical prosodic word. Attributive derivation has its own quantity
conditions and is not automatically shortened by a long vowel in an inner
domain (§\,\ref{sec:stative-attr}).
The distinction between a construction's quantity requirements and its
segmental exponents is compatible with models of morpheme-specific phonology
reviewed in \cite{inkelas2017}; the particular stem classes remain
language-specific.

\section{Denominal verbalisation}\label{sec:denominal-verbalisation}

Among overt noun-to-verb strategies, suffixation in \ird{-ěvat-} and
\ird{-it-} is the most productive. These are the predominant verbalising
formatives with learned, foreign, and recently borrowed nominal bases.
When a foreign stem is adapted directly as a verb, \ird{-ěvat-} is used
exclusively. This construction selects the front grade at the nominal
stem boundary; /cɛ/ and /ɲɛ/ are written \ird{tě} and \ird{ně},
respectively (§\,\ref{sec:morphophon}). Only \ird{-ěvat-} is specified to select a productive front grade here;
the initial \ird{i} in \ird{-it-} does not independently license that mapping. The licensed
boundary is the one between the nominal base and the verbaliser: in
\ird{kontaktěvat-}, PAL applies to the base-final /t/, not to the final
/t/ supplied inside the verbalising exponent. Later inflection establishes
its own boundary on the resulting verbal stem.

\pex
\a \ird{kontakt} $\rightarrow$ \ird{kontaktěvat-} \trsl{contact; get in touch}
\a \ird{telefon} $\rightarrow$ \ird{telefoněvat-} \trsl{telephone}
\xe

The noun \ird{spum} \trsl{bell} supplies the activity stem
\ird{spumit-} \trsl{ring}, in which \ird{-it-} verbalises the nominal
base. Ingressive \ird{ra-} extends lexically to the onset of this activity,
yielding \ird{raspumit-} \trsl{ring out; begin to ring}:

\ex\label{ex:denominal-bell-ingress}%
\begingl
\gla Spum raspúmitek.//
\glb bell.\Dir{} \textsc{ingr}-bell-\textsc{vblz}-\Pf{}//
\glft \trsl{The bell rang out.} / \trsl{The bell began to ring.}
\lingcontext{a contextually identifiable bell; the onset is located in a past narrative}//
\endgl
\xe

The perfective \ird{-ek} presents the onset as a bounded event, while the
activity meaning comes from denominal verbalisation and its ingressive
interpretation from lexical \ird{ra-}
(§\,\ref{sec:lexical-preverbs}). The strong grade lengthens the nominal
base to \ird{spúm-}, leaving verbalising \ird{-it-} and perfective
\ird{-ek} short. The bell is the sole direct-case participant. No past-tense
suffix occurs: the English past renders the narrative context, rather than
an additional meaning of \ird{-it-} or \ird{-ek}
(§\,\ref{sec:vm-asp-case}).

\subsection{Abundance and coverage in \ird{-at-}}\label{sec:abundance-at}

Denominal \ird{-at-} forms a property-denoting stative stem meaning
\trsl{be abundant in, covered with, or characteristically furnished with X}.
Its inputs are semantically compatible substances, materials and kinds whose
presence can characterise a bearer. The construction does not freely turn a
human or institutional noun into a resemblance predicate. Its output belongs
to the stative system (§\,\ref{sec:stative}), rather than to a separate
adjective class, and introduces no additional participant.

The noun \ird{kolen} \trsl{grass} has a lexically specified fleeting
\ird{e}. At this derivational boundary, \ird{-at-} selects its reduced
G0 stem \ird{koln-} (§\,\ref{sec:vowel-zero}):
\ex
\ird{koln-} + \ird{-at-} $\rightarrow$ \ird{kolnat-}
\trsl{be grass-covered; be rich in grass}.
\xe
The \ird{n} belongs to the nominal stem; the formative is \ird{-at-}.
Reduction does not itself select a strong or raised nominal grade, so the
root \ird{o} remains short. The sequence syllabifies /kɔl.nat/, with
ordinary coda lateral vocalisation in
\phon{kolnat-}{ˈkɔw.nat}{grass-covered stative stem}
(§\,\ref{sec:l-vocalisation}). Neither vowel epenthesis nor consonant
loss is required.

The selected G0 nominal base and short formative specify this derived stem,
not a complete inflectional paradigm. The noun's declension and quantitative
class, and the derived stem's finite quantity and ordinary attributive
classes, remain unassigned. Finite aspect and adnominal use require their
respective listed principal parts under §§\,\ref{sec:verbal-quantity}
and~\ref{sec:stative-attr}; a bare stative stem supplies no independent
attributive ending.

Abundance and coverage differ from the relational character or association
expressed by \ird{-(e)zn-} below: being rich in a material need not mean
resembling it or belonging to its social or stylistic domain.
Participant-oriented relative modification (§\,\ref{sec:nm-ovni}) retains
its embedded Voice and aspect selection; it does not contribute the lexical
abundance meaning of \ird{-at-}.

\subsection{Relational statives in \ird{-(e)zn-}}

A semantically specialised but productive denominal pattern adds
\ird{-(e)zn-} to a nominal base to form an intransitive relational stative
verb. The core meaning is \trsl{be X-like; belong to the domain of X; have
the character, style, material, or social colouring of X}. These verbs
conjugate exactly like the ordinary stative class described in
§\,\ref{sec:stative}, and they likewise admit attributive derivation with
its own quantity classes and the manner adverb in \ird{-ě}
(§§\,\ref{sec:stative-attr}--\ref{sec:stative-adv}). Aspect follows the
usual stative pattern: the continuous gives plain state predication, the
progressive requires a licensed stage-level or expressive construal, the
perfective requires a bounded-state construal, and the retrospective describes
an anterior state. None entails entry into the state. Productive
\ird{-él-}/\ird{-el-} instead derives a change-of-state predicate below
viewpoint and alignment (§\,\ref{sec:stative-inchoative}); classificatory
relational meanings require a context supporting an actual transition. The construction selects the restricted
velar-front grade: stem-final \ird{k} surfaces as \ird{c}, as in
\ird{barok} $\rightarrow$ \ird{barocezna}, while coronals remain plain,
as in \ird{kmětezna} and \ird{svícezna}. The independently added manner
suffix \ird{-ě} selects the front grade at the final consonant of the
derived stative stem, yielding \ird{-zně}. The two selections belong to
different boundaries. In \ird{barocezně}, the nominal boundary has already
selected /ts/ from /k/ without coronal PAL, while the adverbial boundary
selects /ɲ/ from the final /n/ of \ird{barocezn-}. Its PAL instruction
does not reach back through \ird{-ezn-} to alter the earlier /ts/.

The choice of segmental \ird{-zn-} or \ird{-ezn-} follows the established
stem class and phonotactics; a long base vowel alone does not delete the
\ird{e}. The formed relational stem supplies a local domain for quantity
selection. Inherited base quantity is retained: long-base relational stems
take short inner aspect endings as Class~C, whereas a short relational stem
without assigned strength follows the suffix-strong pattern. Thus
continuous \ird{tvírzna} has short \ird{-a}, and its progressive and
contemplative take short \ird{-ime} and \ird{-až}. Attributive derivation
is a separate protected layer: \ird{tvírzná} retains both its lexical
\ird{í} and attributive \ird{á}. The same distinction preserves the
listed attributives below rather than imposing a global limit of one long
vowel per word.

Table~\ref{tab:relstatives} lists representative continuous forms together
with their derived attributives and manner adverbs.

\begin{table}[ht]
  \footnotesize
  \caption{Relational stative verbs in \ird{-(e)zn-}.}\label{tab:relstatives}
  \begin{tabularx}{\textwidth}{L l l l L}
    \toprule
    {\sc base noun} & {\sc relational stative verb} & {\sc attributive} & {\sc manner adverb} & {\sc core sense} \\
    \midrule
    \ird{tvír} \trsl{court, courtly sphere} & \ird{tvírzna} & \ird{tvírzná} & \ird{tvírzně} & \trsl{be courtly; belong to the courtly register} \\
    \ird{ščech} \trsl{guild, craft corporation} & \ird{ščehezna} & \ird{ščehezná} & \ird{ščehezně} & \trsl{be guild-like; corporate, artisanal, guild-marked} \\
    \ird{mněra} \trsl{faith, creed} & \ird{mněrazna} & \ird{mněrazná} & \ird{mněrazně} & \trsl{be confessional; marked by a creed or religious orientation} \\
    \ird{kaměn} \trsl{stone} & \ird{kaměnzna} & \ird{kaměnzná} & \ird{kaměnzně} & \trsl{be stony; have a stone-like character; be stone-styled} \\
    \ird{sylo} \trsl{cloth, textile} & \ird{sylozna} & \ird{sylozná} & \ird{sylozně} & \trsl{be textile-like; be clothed, draped, or fabric-marked} \\
    \ird{barok} \trsl{baroque; baroque style} & \ird{barocezna} & \ird{barocezná} & \ird{barocezně} & \trsl{be baroque in style; belong to a baroque aesthetic} \\
    \ird{zmok} \trsl{temple, shrine} & \ird{zmocezna} & \ird{zmocezná} & \ird{zmocezně} & \trsl{be temple-like; be cultic, sacral, or shrine-associated} \\
    \ird{svíc} \trsl{peasant, smallholder, rustic estate} & \ird{svícezna} & \ird{svícezná} & \ird{svícezně} & \trsl{be peasantlike; be rustic or village-class marked} \\
    \ird{kmět} \trsl{law, order, code} & \ird{kmětezna} & \ird{kmětezná} & \ird{kmětezně} & \trsl{be juridical, code-bound, institutionalised} \\
    \bottomrule
  \end{tabularx}
\end{table}

\pex
\a \ird{tvírzná rěč} \trsl{courtly speech; speech of courtly register}
\a \ird{kaměnzná stěna} \trsl{a wall with a stony character} \lingcontext{a literal \trsl{stone wall} is expressed analytically, not by \ird{kaměnzná}}
\a \ird{barocezná dvorba} \trsl{a baroque-style building}
\xe


Coverage and relational character give different grounds for a property
predication. The abundance derivative in (\ref{ex:deriv-slot-coverage}a)
uses the stative stem \ird{kolnat-}; its finite quantity class is left
unspecified by the complete-cell notation:

\pex\label{ex:deriv-slot-coverage}
\a
\begingl
\gla \textnormal{[[garden]]}\textsubscript{\Dir{}} \textnormal{[[grass-covered]]}\textsubscript{\Cont{}}.//
\glb garden.\Dir{} grass-\textsc{abund}-\Cont{}//
\glft \trsl{The garden is covered with grass.}
\lingcontext{grass grows over the garden; the predication concerns coverage}//
\endgl
\a
\begingl
\gla \textnormal{[[wall]]}\textsubscript{\Dir{}} kaměnzna.//
\glb wall.\Dir{} stone-\textsc{rel}-\Cont{}//
\glft \trsl{The wall looks like stone.}
\lingcontext{paint gives the wall the appearance of stone; the wall need not be made of stone}//
\endgl
\xe

The relational predicate in (b) can characterise an appearance without
asserting the presence of stone as a covering material. Conversely, actual
grass coverage in (a) supplies no claim that the garden resembles grass
or belongs to a grass-associated stylistic domain.

\section{Lexical preverbs}\label{sec:lexical-preverbs}

The lexical preverbs specify a path, a configuration, or a relation to an
earlier outcome within the meaning of a derived verb. Directional derivation
includes goal-oriented \ird{de-}, outward \ird{le-}, ingressive \ird{ra-},
circuitous \ird{vra-}, and traversal \ird{mně-}; configurational
\ird{ša-} and \ird{na-} express dispersal and accumulation, respectively.
Restitutive \ird{ru-} relates a path or state to an earlier configuration,
whereas \ird{pře-} supplies replacement or revision. Reversative \ird{po-}
cancels a resultant condition or configuration. Marginal privative
\ird{ne-} remains confined to individually specified derivatives. These
semantic groupings describe the lexical contributions of the formatives;
they introduce no additional morphological positions.
Table~\ref{tab:lexical-preverbs} gives their principal distributions, with
hyphens marking bound formatives and lexical stems.

\begin{table}[htbp]
  \footnotesize
  \caption{Lexical preverbs and their principal distributions.}\label{tab:lexical-preverbs}
  \begin{tabularx}{\textwidth}{l L L}
    \toprule
    {\sc preverb} & {\sc contribution and distribution} & {\sc derivative} \\
    \midrule
    \ird{de-} & Goal or endpoint orientation; productive with compatible motion predicates
      & \ird{denočk-} \trsl{run to a destination} \\
    \ird{le-} & Emergence from a bounded domain; productive with source or containment predicates
      & \ird{lenočk-} \trsl{run out} \\
    \ird{ra-} & Productive ingress with compatible entry predicates; lexical extensions to activity onset
      & \ird{ranočk-} \trsl{run into an enclosure}; \ird{raspumit-} \trsl{begin to ring} \\
    \ird{vra-} & Circumvention or indirect course; moderately productive with compatible paths
      & \ird{vranočk-} \trsl{run around an obstacle} \\
    \ird{mně-} & Traversal of a path or ordered domain; productive with coherent courses
      & \ird{mněkař-} \trsl{read through}; \ird{mněvad-} \trsl{think through} \\
    \ird{ša-} & Divergence from a common configuration; restricted to dispersive predicates
      & \ird{šastoj-} \trsl{go off in different directions} \\
    \ird{na-} & Accumulation or assembly; productive with compatible configurations
      & \ird{navad-} \trsl{synthesise considerations into a judgement} \\
    \ird{ru-} & Retracing or restitution; productive with recoverable earlier configurations
      & \ird{rustoj-} \trsl{retrace an earlier route} \\
    \ird{pře-} & Replacement or revision; restricted chiefly to cognition, assessment, and transformation
      & \ird{převad-} \trsl{reconsider; revise a judgement} \\
    \ird{po-} & Undoing, removal, or dismantling; compatible resultant conditions and configurations
      & \ird{poshod-} \trsl{take off clothing} \\
    \ird{ne-} & Privation, removal, or dispossession; individually specified derivatives
      & No general \trsl{undo} or negative rule \\
    \bottomrule
  \end{tabularx}
\end{table}

\subsection{Lexical families and semantic selection}

The motion family based on \ird{nočk-} \trsl{run} distinguishes the course
of an event without changing its manner. Unprefixed \ird{nočk-} specifies
no inherent direction, while \ird{denočk-} orients the running towards a
destination. The stem \ird{lenočk-} specifies emergence from an enclosed or
bounded region, \ird{ranočk-} entry into such a region, and \ird{vranočk-} a
circuitous course around an obstacle or away from a direct route. Orientation
towards a destination in \ird{de-} does not itself specify the boundary
crossing contributed by \ird{ra-}; similarly, \ird{le-} requires a
recoverable source domain and does not apply to departure irrespective of
containment. Concealment and emergence can support extensions of \ird{le-}
to extraction or disclosure, just as entry into a conceptual domain can support
\ird{ra-} with a non-spatial predicate. Such extensions depend on the meaning
of a particular derivative rather than on unrestricted combination with every
communication or cognition verb.
Activity onset provides a further lexical extension: in
\ird{raspumit-} \trsl{begin to ring}, ingress concerns entry into the
ringing activity derived from \ird{spum} \trsl{bell}
(\ref{ex:denominal-bell-ingress}). The onset interpretation belongs to this
prefixed stem, whose viewpoint is supplied separately by aspect morphology.

The preverb \ird{vra-} requires a possible direct course against which
circumvention or indirectness can be understood. Its spatial meaning can therefore motivate
lexical extensions to indirect inquiry or strategic avoidance, although the
mere presence of an extended process does not license it. In
\ird{mněkař-}, by contrast, the text supplies an ordered domain through
which reading proceeds, while \ird{mněvad-} extends the same organisation
to the stages of considering a problem. The preverb \ird{mně-} thus
concerns traversal through an ordered domain, whereas \ird{vra-} concerns the course taken
around or in avoidance of one.

With \ird{stoj-} \trsl{go; directed motion}, \ird{šastoj-} describes
movement along diverging trajectories from a common origin or configuration.
The context must support more than one trajectory, normally through multiple
participants; the preverb neither supplies grammatical plural marking nor
licenses an otherwise unexpressed participant. Dispersal concerns the spatial
or processual configuration of the event, rather than intermittent action,
repeated small events, or scattered iteration, and supplies no reduplicative
operation. The restriction to genuinely diverging configurations accounts for the limited
productivity of \ird{ša-}.

The preverb \ird{na-} supplies the converse configurational relation: entities
or processes accumulate or are assembled into a common configuration. The productive spatial
pattern extends to gathering and integration where the base already admits
such a relation. In the cognitive family \ird{vad-} \trsl{think},
\ird{mněvad-} \trsl{think through a problem}, and \ird{převad-}
\trsl{reconsider or revise a judgement}, conventionalised \ird{navad-}
\trsl{bring separate considerations into a coherent judgement} supplies a
further distinction. Its synthesis meaning is lexically specified and does not
extend mechanically to every cognition verb. The considerations and judgement
belong to that lexical interpretation; prefixation does not thereby license a
new nominal object. Mere accompaniment or collective agency supplies no
accumulative meaning, and an added applied argument requires the ordinary
applicative construction.

The stem \ird{rustoj-} \trsl{retrace an earlier route} relates the current
movement to a previously established path. More generally, \ird{ru-} requires an
earlier configuration towards which the event returns, whether this is a route
or a restorable condition. The preverb \ird{pře-} instead concerns
replacement or revision of an earlier outcome, as in \ird{převad-}; a revised judgement need not
restore an earlier one. Renewed occurrence at discourse level, including the
reading of \ird{jas vy} (§\,\ref{sec:class2-particles}), likewise need not
involve lexical restitution: that construction requires earlier validity
and an intervening contrary stage without specifying an earlier path or a
sublexical result-state attachment site. Privative \ird{ne-} contributes removal or
deprivation wherever a derivative licenses it, without expressing the failure
or non-occurrence of the base event.


A destination can guide a moving participant without an enclosure being
entered. The following pair holds the mover and progressive viewpoint
constant while changing that path specification:

\pex\label{ex:deriv-slot-path}
\a
\begingl
\gla \textnormal{[[dog]]}\textsubscript{\Dir{}} denóčkime.//
\glb dog.\Dir{} \textsc{goal}-run-\Prog{}//
\glft \trsl{The dog is running towards the house.}
\lingcontext{the house is the destination supplied by context; entry is not implied}//
\endgl
\a
\begingl
\gla \textnormal{[[dog]]}\textsubscript{\Dir{}} ranóčkime.//
\glb dog.\Dir{} \textsc{ingr}-run-\Prog{}//
\glft \trsl{The dog is running into the house.}
\lingcontext{the same house supplies the enclosure; the running follows a course through its entrance}//
\endgl
\xe

Neither progressive asserts completed arrival or entry. The enclosure
belongs to the interpretation of the path in (b), without becoming an
additional direct object; an overt spatial dependent would retain the
ordinary postpositional construction.

\subsection{Reversal, removal, and dismantling in \ird{po-}}\label{sec:reversative-po}

Lexical \ird{po-} selects a base whose interpretation supplies a recoverable
resultant condition, possession, or established connection. Its productive
core includes reversal of that condition, privative removal of material,
an object or a property, and deconfigurative dissolution of an arrangement.
These readings cancel or dismantle what is established; they need not restore
an exact earlier state. Oppositional meanings require individual lexical
specification, rather than a productive rule negating any state or action.

The existing \ird{shod-} \trsl{wear; put on} supplies
\ird{poshod-} \trsl{take off clothing}, reversing its putting-on sense.
The clothing relation is retained; prefixation supplies no new core argument.
The Class~A progressive
\phon{poshódime}{ˈpɔ.sxuː.dɪ.mɛ}{is taking off clothing}
selects root G1 while stress falls on \ird{po-}; it need not describe
completed removal. Linked \ird{poshódimně} retains the same root quantity.
Indicative negation is external:
\phon{zaposhódime}{zaˈpɔ.sxuː.dɪ.mɛ}{is not taking off clothing}.
It denies the removal event rather than supplying an inverse of the base.

Restitutive \ird{ru-} returns towards an earlier configuration, whereas
\ird{po-} removes or dismantles a resultant one. Revising \ird{pře-}
can supply a replacement outcome; \ird{po-} need not supply one.
Marginal privative \ird{ne-} retains its individually specified distribution;
it is not a productive allomorph of \ird{po-}.
Like the other lexical preverbs, \ird{po-} precedes ordinary valency and
aspect morphology, selects no independent mutation or quantity grade, and
supplies neither perfectivity nor an automatic inverse of every verb.
Productive stacking with another lexical preverb is unavailable.

This lexical preverb is distinct from homophonous stative degree \ird{po-}
\trsl{slightly}, as in \ird{pomužna} \trsl{is somewhat beautiful}
(§\,\ref{sec:stative-degree}), and from the free complementizer \ird{po}
(Chapter~\ref{ch:embedding}). No productive combination of the lexical and
degree prefixes is specified here. Existing lexical stems are not segmented
merely because their first syllable is \ird{po}.


Removing a garment and denying its removal differ even when the same
participants and viewpoint are retained:

\pex\label{ex:deriv-slot-reversal}
\a
\begingl
\gla Veňo hráťam poshódime.//
\glb woman.\Dir{}.\Def{} coat.\Ind{} \textsc{rev}-wear-\Prog{}//
\glft \trsl{The woman is taking off a coat.}
\lingcontext{removal is under way; the coat need not yet be off}//
\endgl
\a
\begingl
\gla Veňo hráťam zaposhódime.//
\glb woman.\Dir{}.\Def{} coat.\Ind{} \Neg{}-\textsc{rev}-wear-\Prog{}//
\glft \trsl{The woman is not taking off a coat.}
\lingcontext{the removal is denied; putting the coat on is not asserted}//
\endgl
\xe

Lexical reversal remains inside the predicate denied in (b). Both clauses
retain the progressive agent-direct, garment-indirect frame, so the
contrast concerns negation of removal rather than a change in participant
selection.

\subsection{Stem formation, aspect, and argument structure}
\label{sec:preverb-aspect}

A preverb and its base form one lexical constituent, to which the ordinary
verbal morphology of §\,\ref{sec:morph-template} subsequently attaches. The
schematic organisation is therefore:

\begin{center}
\small
\begin{tabular}{c}
(\Neg{}) + [preverb + root]$_{\text{stem}}$ + (Antip/Appl)\\
+ (Caus) + (Voice \ird{-n-}/\ird{-an-}) + Asp.Align + (Mood) + (Evid)
\end{tabular}
\end{center}

The bracketed constituent occupies the stem position of the finite template;
prefixation supplies no additional inflectional category. In the analysis
adopted here, the lexical constituent is interpreted within the predicate,
below the functional heads which license participants, causation, and viewpoint
aspect. The suffixal hierarchy described in Chapter~\ref{ch:argument-structure} consequently
applies to prefixed and unprefixed stems alike. A derived stem retains its
base's thematic class and construction-selected grades; the presence of a
preverb does not replace the ordinary Voice, antipassive, applicative, or
causative operations.

The distinction between lexical event structure and grammatical viewpoint has
a close comparative parallel in Hungarian, where a telic predicate can receive
a progressive interpretation \cite{kardosfarkas2022}. Hungarian
separable particles also alternate in position under focus, negation, and
progressive aspect \cite{kardosfarkas2022}; Iridian preverbs remain
bound throughout these environments. The comparison concerns the independence
of the two aspectual dimensions, without extending the Hungarian syntactic
analysis to Iridian.

All eleven preverbs remain compatible with continuous, progressive, contemplative,
perfective, and retrospective morphology wherever the resulting predicate
admits the relevant interpretation. None supplies grammatical perfectivity.
Thus progressive \ird{lenóčkime} \trsl{is running out},
\ird{ranóčkime} \trsl{is running into an enclosure}, and
\ird{rustójime} \trsl{is retracing a route} describe the unfolding of the
specified course without asserting its completion. These Class~A forms select
G1 in the root and short \ird{-ime}, as do
\ird{vranóčkime}, \ird{šastójime}, and \ird{navádime}.

The aspectual paradigm of \ird{mněkař-} illustrates the independence of
lexical traversal from grammatical viewpoint. Reading through a text has a
potential endpoint, but neither the continuous nor the progressive asserts
that the endpoint has been reached. The contemplative projects the reading
event, whereas the perfective presents it as a bounded whole and the
retrospective relates the prior event to a subsequent vantage point. The
transitive forms in Table~\ref{tab:preverb-aspect} retain the text as patient;
the perfective and retrospective therefore include overt Voice \ird{-n-}.

\begin{table}[htbp]
  \footnotesize
  \caption{Transitive \ird{mněkař-} across the five aspects.}\label{tab:preverb-aspect}
  \begin{tabularx}{\textwidth}{l l L}
    \toprule
    {\sc aspect} & {\sc finite form} & {\sc interpretation} \\
    \midrule
    Continuous & \ird{mněkařa} & Be engaged in reading through a text; completion is not asserted \\
    Progressive & \ird{mněkářime} & Be actively proceeding through a text \\
    Contemplative & \ird{mněkařáž} & Intend or expect to read through a text \\
    Perfective & \ird{mněkářnik} & Read through a text as a bounded completed event \\
    Retrospective & \ird{mněkařnáv} & Have read through a text, viewed from a subsequent reference point \\
    \bottomrule
  \end{tabularx}
\end{table}

Prefixation itself neither licenses an additional argument nor changes case
assignment. Thus \ird{denočk-} retains the intransitive motion frame of
\ird{nočk-}, with any overt destination expressed by the ordinary goal
construction, while \ird{mněkař-} retains the transitive frame of
\ird{kař-}. In the perfective and retrospective domains, an ordinary
transitive derivative takes Voice \ird{-n-} and ergative alignment; its
antipassive instead follows the independently established patient-suppressing
construction. The continuous, progressive, and contemplative retain nominative
alignment (§\,\ref{sec:vm-asp-case}). The motion derivatives \ird{lenočk-},
\ird{ranočk-},
\ird{vranočk-}, \ird{rustoj-}, and \ird{šastoj-} likewise retain
an intransitive mover frame: a source, enclosure, obstacle, or earlier route
is not made a direct object by prefixation. Overt spatial dependents follow
the ordinary postpositional constructions (§\,\ref{sec:nm-postpositions}).
The stem \ird{navad-} preserves the cognitive frame of \ird{vad-}. An exceptional
lexical case frame or stem alternant requires specification with the particular
derivative; it does not follow from the presence of a preverb.

The unprefixed \ird{kař-} can itself occur in the perfective, which bounds
a reading event without adding the traversal meaning of \ird{mně-}.
Conversely, the continuous and progressive already provide internal viewpoints
on \ird{mněkař-}; prefixed stems require no separate imperfectivising
derivation. The phasal particles \ird{vy} and \ird{so} retain their ordinary
scope over the predicate and contribute the interpretations described in
§\,\ref{sec:class1-particles}, independently of the lexical preverb.

\subsection{Bound status and lexical restrictions}

A lexical preverb remains adjacent to its base under focus, negation,
interrogation, and changes of mood. It is written together with the inflected
stem, and no nominal, particle, or other syntactic constituent can intervene at the preverb--root
boundary. Indicative \ird{za-} attaches outside the derived stem, as in
\ird{zamněkářime} \trsl{is not reading through}; the lexical privative
\ird{ne-} occupies a different position and supplies no substitute for
indicative negation (§\,\ref{sec:vm-neg}). Outward motion is supplied by
\ird{le-}, as in \ird{lenóčkime}; \ird{zalenóčkime}
\trsl{is not running out} places negative \ird{za-} outside that stem,
without assigning it a lexical directional meaning. Productive derivation
admits one lexical preverb per base. An apparently complex borrowed or lexicalised stem,
such as \ird{prit-} or \ird{podohlet-}, is not thereby segmented into a
productive native prefix and root.

The preverb belongs to the minimal lexical prosodic word and receives initial
primary stress: \ird{mněkářime} and \ird{zamněkářime} are both stressed
on \ird{mně-}, while \ird{převádime} is stressed on \ird{pře-}
(§\,\ref{sec:stress}). The four lexical syllables of \ird{mněkářime}
receive subsidiary prominence on the third syllable, \ird{ři-}; external
\ird{za-} adds no syllable to that calculation, so \ird{zamněkářime}
retains the same footing (§\,\ref{sec:secondary-stress}). Prefixation
introduces no independent mutation or
quantity grade. The derivative follows the ordinary construction-selected
stem grades, suffixal consonant mutations, and cluster phonology, with any
lexically exceptional allomorph recorded separately. Each derived stem has its
own lexical entry, specifying its base, meaning, argument structure, and any
special alternants. The consonant-initial bases illustrated here establish no
productive rule for a vowel-initial root: combinations at that boundary require
lexical specification rather than automatic epenthesis, contraction, deletion,
or glide insertion (Chapter~\ref{ch:phon}).

Quantity selection excludes the preverb from its root target. Thus the
progressive \ird{mněkářime} lengthens the root nucleus while primary stress
stays on \ird{mně-}. Contemplative \ird{mněkařáž} and retrospective
\ird{mněkařnáv} instead retain the plain root and their long aspect
endings. Root quantity, suffix quantity, and stress are independently specified.

\section{Property and stative derivation}\label{sec:stative-derivation}

\subsection{Property roots and the verbalising suffix \ird{-n-}}

Property roots feed productive stative derivation in \ird{-n-}; their
inventory and their attributive class remain lexically specified. The
stativiser supplies no extra participant. Lexical positional, possessive and
locative states need not bear \ird{-n-}: semantic stativity and this particular
morphological class are separate properties.

\begin{table}[ht]
\footnotesize
\caption{Property roots and derived stative verbs.}\label{tab:statroots}
\begin{tabularx}{\textwidth}{l L l L}
\toprule
{\sc root} & {\sc gloss} & {\sc stative stem} & {\sc gloss} \\
\midrule
\ird{hrad} & goodness & \ird{hradn-} & be good \\
\ird{muž} & beauty & \ird{mužn-} & be beautiful \\
\ird{torba} & badness & \ird{torban-} & be bad \\
\ird{svor} & bigness & \ird{svorn-} & be big \\
\bottomrule
\end{tabularx}
\end{table}

In \ird{mužna} /mʊʒna/, the voiced-obstruent--nasal sequence /ʒn/
admits an onset under the rising-sonority pattern of §\,\ref{sec:sonority}.
Stativisation therefore retains both root-final \ird{ž} and \ird{-n-};
no epenthesis, deletion or additional palatalisation is required. Strong
root grade is \ird{múžn-}, with /uː/, whereas plain linking gives
\ird{mužne} and adverbial PAL gives \ird{mužně}.

\begin{table}[htbp]
\footnotesize
\caption{Constructed state and inchoative principal parts of \ird{mužn-}.}
\label{tab:stative-become}
\begin{tabularx}{\textwidth}{l L L}
\toprule
Aspect & Simple state & Overt inchoative \\
\midrule
Continuous & \ird{mužna}: is beautiful & \ird{mužnéla}: becomes beautiful \\
Progressive & \ird{múžnime}: looks particularly beautiful & \ird{mužnélime}: is becoming beautiful \\
Perfective & \ird{múžnik}: was beautiful for the bounded occasion & \ird{mužnélek}: underwent becoming beautiful \\
Retrospective & \ird{mužnáv}: had been beautiful & \ird{mužnélav}: had become beautiful \\
Contemplative & \ird{mužnáž}: is expected to be beautiful & \ird{mužnélaž}: is expected to become beautiful \\
\bottomrule
\end{tabularx}
\end{table}

The simple progressive requires an expressive or stage-level context; the
simple perfective requires a context that licenses a bounded state. Neither
is a free substitute for a transition, and an enduring quality without such
context does not license those readings. Retrospective state predication is
anterior and does not entail a preceding change or cessation of the state.
Contemplative is projective, not a future tense, and does not entail that
the projected state has not already begun. These constructed paradigms specify grammatical contrasts,
not claims of native attestation.

The finite principal parts of \ird{mužn-}, \ird{hradn-},
\ird{svorn-}, \ird{četn-}, \ird{ptakn-}, \ird{všihn-},
\ird{hvodn-} and \ird{žbirn-} belong to root-strong Class~A: G0 continuous, retrospective
and contemplative, G1 progressive and perfective. These specified finite
classes do not determine their separate attributive classes. Protected
\ird{tvírzn-} belongs to Class~C and retains its lexical quantity with
short outer aspect endings. Short relational stems retain their independently
listed suffix-strong pattern; the semantic classes of §\,\ref{sec:stative-stability} do not overwrite it.
The recorded entry \ird{smrn-} lacks a specified nucleus under the
general phonology. Its spelling alone licenses neither a repaired stem nor
additional finite or attributive outputs.

The property root \ird{torba} \trsl{badness} supplies stative
\ird{torban-} \trsl{be bad}, with syllabification /tɔr.ban/.
Its recorded quantity target is the final root \ird{a}, excluding
stativising \ird{-n-} and outer suffixes: G0 \ird{torban-} contrasts
with conditional G1 \ird{torbán-}. The initial \ird{o} remains short,
and primary stress remains initial. Finite class, Become class and ordinary
attributive class are unassigned for this lexeme. The following outputs
are conditional paradigms under the specified assignments, not licensed
principal parts; no ordinary attributive or manner form is selected by
its phonological shape alone.

\begin{table}[htbp]
\footnotesize
\caption{Conditional \ird{torban-} paradigms: finite Class~A and Become EL-S.}
\label{tab:torban-conditional}
\begin{tabularx}{\textwidth}{l L L L L}
\toprule
Aspect & State & Linked state & Become & Linked Become \\
\midrule
Continuous & \ird{torbana} & \ird{torbane} & \ird{torbanéla} & \ird{torbanéle} \\
Progressive & \ird{torbánime} & \ird{torbánimně} & \ird{torbanélime} & \ird{torbanélimně} \\
Perfective & \ird{torbánik} & \ird{torbánice} & \ird{torbanélek} & \ird{torbanélice} \\
Retrospective & \ird{torbanáv} & \ird{torbanávě} & \ird{torbanélav} & \ird{torbanélavě} \\
Contemplative & \ird{torbanáž} & \ird{torbanáže} & \ird{torbanélaž} & \ird{torbanélaže} \\
\bottomrule
\end{tabularx}
\end{table}

Simple-state progressive and perfective retain their stage-level and
bounded-state conditions. In the conditional EL-S paradigm, Become selects
local \ird{-él-}; outer aspect does not lengthen either property-root
vowel through that head. The manner candidate \ird{torbaně} requires
its own lexical licence. Ordinary attributive class remains unlisted;
none of the three ordinary ATTR classes is implied by finite Class~A.

The stative verbaliser selects perfective \ird{-ik}, as in root-strong
\ird{múžnik} and \ird{svórnik}. The suffix does not add transitive
orientation. A stative applicative retains its independently specified
\ird{-ik}, while the intervening valency head supplies the local quantity
rather than permitting a new aspectual root grade.

\subsection{Overt inchoative derivation}\label{sec:stative-inchoative}

Productive Become takes a stative predicate and derives a single-participant
transition predicate. Its exponent is plain \ird{-el-}, strong \ird{-él-}.
Quantitative selection belongs to a separately listed derivational class.
The strong extension class contains \ird{mužn-}, \ird{hradn-},
\ird{svorn-}, \ird{četn-}, \ird{ptakn-}, \ird{všihn-},
\ird{hvodn-} and \ird{žbirn-}, whose base-grade
inputs give \ird{mužnél-}, \ird{hradnél-}, \ird{svornél-},
\ird{četnél-}, \ird{ptaknél-}, \ird{všihnél-}, \ird{hvodnél-}
and \ird{žbirnél-}. The protected-input class has
plain \ird{-el-}, as in listed \ird{tvírznel-}. A short root alone
does not assign an unlisted predicate to the strong extension class.
The unprotected short input is its base grade, not the grade which an outer
progressive might otherwise have selected in a simple state.

Become establishes a local quantity-bearing derived stem. Its adjacent outer
aspect exponents are short \ird{-a}, \ird{-ime}, free perfective
\ird{-ek}, \ird{-av} and \ird{-až}. Inchoative perfective is eventive:
\ird{mužnélek}, not a mechanical extension of the stative \ird{-ik}.
The ordinary bound high grade gives participant \ird{mužnélikou} and
relative \ird{mužnélikový}; front-grade linking gives \ird{mužnélice}.
Asp cannot reach through Become to assign a fresh grade to the property root.
An intervening valency head establishes its own domain under the existing
rules. Become precedes inner Appl, Caus and outer orientation; it does not
interchange their positions or itself license a causer, stimulus or patient.

The derivation asserts change along a lexically or contextually supplied
scale. The distinction between scalar motion and variable telicity has a
comparative parallel in the analysis of Korean \textit{-eci}
\cite{limzubizarreta2012}. An open scale admits increase without a fixed endpoint; a threshold or
closed endpoint must be independently supplied. Perfective bounds the change
but does not by itself assert attainment of an arbitrary positive standard.
Thus \ird{svorna} is \trsl{is big}; \ird{svornéla} can mean
\trsl{becomes bigger} with a comparative change standard, or
\trsl{becomes big} when a positive threshold is supplied. Likewise
\ird{svornéra} means \trsl{makes bigger} under a comparative standard and
\trsl{makes big} under an independently specified positive threshold.
Causation, scalar increase and threshold attainment are distinct contributions.
These five interpretations consequently involve three morphological stems:
the comparative versus positive standard must be supplied by an overt
comparison construction or an independently recoverable scalar context
(Chapter~\ref{ch:predicates}). The endings alone do not disambiguate the two
change standards; no additional threshold suffix is licensed.


\begin{table}[htbp]
\footnotesize
\caption{Constructed scalar contrasts with independently specified standards.}
\label{tab:stative-scalar-contrasts}
\begin{tabularx}{\textwidth}{l L L}
\toprule
Form & Interpretation & Scalar specification \\
\midrule
\ird{svorna} & be big & state degree meets a positive standard \\
\ird{svornéla} & become bigger & increase relative to an earlier degree; no positive threshold required \\
\ird{svornéla} & become big & transition whose target meets the supplied positive threshold \\
\ird{svornéra} & make bigger & caused increase; no positive threshold required \\
\ird{svornéra} & make big & caused transition whose target meets the positive threshold \\
\bottomrule
\end{tabularx}
\end{table}

These are semantic contrasts with supplied standards, not five distinct
exponents. A transition's target need not be attained in an ongoing
continuous or progressive description. An overt comparative with
\ird{stá} supplies a comparative target property, such as exceeding a
named standard; it does not by itself entail mere increase relative to the
holder's own earlier degree (§\,\ref{sec:comparison}).


A child can grow without yet being big for their age; the following
contrast distinguishes increase in size from a positive size standard:

\pex\label{ex:deriv-slot-scalar}
\a
\begingl
\gla Byl svorna.//
\glb child.\Dir{} big-\textsc{state}-\Cont{}//
\glft \trsl{The child is big.}
\lingcontext{the child is big for their age; no change is asserted}//
\endgl
\a
\begingl
\gla Byl svornéla.//
\glb child.\Dir{} big-\textsc{state}-\textsc{become}-\Cont{}//
\glft \trsl{The child is getting bigger.}
\lingcontext{the child has grown but can still be small for their age}//
\endgl
\a
\begingl
\gla Byl svornéla.//
\glb child.\Dir{} big-\textsc{state}-\textsc{become}-\Cont{}//
\glft \trsl{The child is getting big.}
\lingcontext{the context supplies the positive standard as the transition's target; attainment is not asserted}//
\endgl
\xe

The identical change forms in (b--c) leave the choice of standard to the
specified contexts. The simple state in (a) entails no preceding growth,
while neither ongoing change description entails that its target has been
reached.

\subsection{Restricted manifestative derivation}\label{sec:stative-manifest}

Manifest derives an outward behavioural or expressive predicate from a
lexically selected internal-condition predicate. The initial licensed input
is \ird{všihn-} \trsl{be happy}, with plain \ird{-eš-}:
\ird{všihneša} \trsl{displays happiness}. The display may be feigned;
the construction does not entail that its bearer genuinely feels happy.
The selected StateP supplies the description of the internal condition,
without asserting that an actual state satisfying that description obtains.
It is consequently distinct both from an assertion of the state and from
indirectly evidenced \ird{všihnal}. It reports neither an evidential
source nor a causal relation, and introduces no external causer.

The derived stem has its own Class~A quantitative domain, whose eligible
nucleus is the \ird{e} of Manifest. Its G1 grade is \ird{é}; an outer
selector cannot reach through Manifest to lengthen the \ird{i} in
\ird{všihn-}. The following principal parts therefore preserve the inner
state stem while distinguishing local derivational quantity from stress:

\begin{table}[htbp]
\footnotesize
\caption{Constructed principal parts of licensed \ird{všihneš-}.}
\label{tab:stative-manifest}
\begin{tabularx}{\textwidth}{l l L}
\toprule
Aspect & Form & Interpretation \\
\midrule
Continuous & \ird{všihneša} & displays happiness \\
Progressive & \ird{všihnéšime} & is displaying happiness as an unfolding activity \\
Perfective & \ird{všihnéšek} & displayed happiness on a bounded occasion \\
Retrospective & \ird{všihnešáv} & had displayed happiness \\
Contemplative & \ird{všihnešáž} & is projected to display happiness \\
\bottomrule
\end{tabularx}
\end{table}

The ordinary eventive perfective has free \ird{-ek} and bound
\ird{-ik-}: participant \ird{všihnéšikou} denotes the selected display
bearer, not happiness or the display event. Its direct linked form is
\ird{všihnéšice}; indirect Evid instead follows the free finite host in
\ird{všihnéšekal}, and further linking yields \ird{všihnéšekalě}.
Plain Manifest attachment retains \ird{n}; the \ird{š} belongs to the
new head and is not a mutation of that nasal. In broad
\phon{všihneša}{ˈfʃɪx.nɛ.ʃa}{displays happiness}, all three nuclei are
short and initial stress is preserved. Progressive
\phon{všihnéšime}{ˈfʃɪx.niː.ʃɪ.mɛ}{is displaying happiness}
selects local G1 /iː/, written \ird{é}, without moving stress;
conservative speech retains /eː/ in this grade.

ManifestP selects StateP below causative and viewpoint structure. In this
restricted paradigm it preserves one display bearer and licenses no stimulus
applicative by inheritance from \ird{všihn-}. Become and Manifest are not
freely stacked, and dimensional, relational and ordinary bodily-condition
predicates do not acquire Manifest simply by having stative meanings.
An independently causable display can take ordinary higher Caus: compact
\ird{všihnéšrime} assigns G1 to the accessible Manifest nucleus, not
the inner root. It means causing a display and need not mean causing genuine
happiness; its two-participant frame follows the ordinary causative rules.

The exponent is distinct from emphatic ATTR \ird{-ěsk-}, which selects
its own front boundary and contains final \ird{k}. Formal resemblance between
parts of different exponents does not establish a productive segmentation
or enlarge the distribution of Manifest or emphatic attribution.


A child may look happy without feeling happy, whereas a report can
describe the feeling without describing any visible behaviour:

\pex\label{ex:deriv-slot-display}
\a
\begingl
\gla Byl všihneša.//
\glb child.\Dir{} happy-\textsc{manifest}-\Cont{}//
\glft \trsl{The child looks happy.}
\lingcontext{the child smiles and laughs; the display need not reflect their feelings}//
\endgl
\a
\begingl
\gla Byl všihnal.//
\glb child.\Dir{} happy-\Cont{}-\textsc{ind.evid}//
\glft \trsl{The child is reportedly happy.}
\lingcontext{someone has reported the child's feelings; no outward display is described}//
\endgl
\xe

The report in (b) concerns the internal state, whereas (a) describes
behaviour whose sincerity remains open. Manifest supplies the display
predicate; Evid supplies the source of knowledge for the state assertion.

\subsection{Attributive and adverbial derivatives}\label{sec:stative-attr}

Ordinary ATTR selects a property-denoting base or reduced StateP, not an
independently finite verb from which aspect is subsequently removed:
\[
[\mathrm{AttrP}\ [\mathrm{PropertyP}/\mathrm{StateP}\ \ldots]\ \mathrm{ATTR}].
\]
The category-neutral root contributes the property, the stativising layer
licenses its holder in a state description, and ATTR forms an adnominal
property without obligatory finite Asp.Align. Surface reuse of the stem in
\ird{-n-} therefore does not establish a hidden finite clause. Its three
listed morphological classes are root-strong, suffix-strong and soft. Root-strong \ird{múžna} selects G1 on the
property root and short \ird{-a}, while suffix-strong \ird{hradná},
\ird{svorná}, and \ird{skorná} retain their base grade with long
\ird{-á}. The attributive forms a separate derivational domain: protected
\ird{tvírzná} can retain both the lexical root quantity and long
attributive ending. These are class-selected correspondences, not a general
shortening of every ending after a long root. The construction is tenseless,
carries no aspectual or alignment marking, and is used exclusively as a
direct modifier within the noun phrase. Attributive class supplies no
nominal definiteness, case agreement or temporal anchor. The lexical meaning
and context determine the property reading; classes do not encode a universal
dispositional-versus-episodic contrast:

\pex
\a \ird{múžna nuž} \trsl{a beautiful city}
\a \ird{hradná byl} \trsl{a good child}
\a \ird{svorná dum} \trsl{a big house}
\xe

The attributive form cannot be used predicatively. For predication, the full
stative verb form is required:

\pex
\a \ird{Nužo mužna.} \trsl{The city is beautiful.} \lingcontext{predicative: stative verb required}
\a \ird{*Nužo múžna.} \lingcontext{ungrammatical: attributive form in predicative position}
\xe

The soft Class~III is listed for \ird{všihn-} and \ird{žbirn-}.
It selects the existing local front grade at the stative boundary, so final
\ird{n} becomes \ird{ň}, and adds short \ird{-a}: \ird{všihňa},
\ird{žbirňa}. The root remains G0; neither semantic membership in an
experiential class nor a short root licenses this formation on another stem.
These forms retain the lexical property meanings \trsl{happy} and
\trsl{wet}, respectively. The following contrasts have different selected
endings, without any new global palatalisation rule:

\pex
\a \phon{všihňa}{ˈfʃɪx.ɲa}{happy, attributive}
\a \phon{všihně}{ˈfʃɪx.ɲɛ}{happily, manner}
\a \phon{žbirňa}{ˈʒbɪr.ɲa}{wet, attributive}
\a \phon{žbirně}{ˈʒbɪr.ɲɛ}{wetly, manner}
\xe

\noindent The attributive functions as the adnominal counterpart of
the continuous stative, expressing the property directly without tense or
aspect. Its position within the noun phrase relative to determiners and other
modifiers is treated in Chapter~\ref{ch:nouns}.

Alongside these three attributive classes, property-root-based stative stems
productively form a dedicated class of derived manner adverbs in \ird{-ě}
\label{sec:stative-adv}. This formation is derivational rather than
inflectional: it yields a lexical adverbial adjunct, not a further verbal
form, and it is distinct both from the adnominal attributive and from the
converb system discussed in Chapter~\ref{ch:complex}. Schematically, the
pattern is as follows:

\begin{center}
\term{property root $\rightarrow$ stative stem in \ird{-n-} $\rightarrow$ class-selected attributive / manner adverb in \ird{-ě}}
\end{center}

\noindent A representative paradigm is the following:

\pex
\a \ird{skor} \trsl{quickness}
\a \ird{skorn-} \trsl{be quick}
\a \ird{skorná} \trsl{quick, fast} \lingcontext{attributive}
\a \ird{skorně} \trsl{quickly}
\xe

The resulting system preserves a clear three-way distinction: the finite stative
predicate expresses property predication, the attributive modifies a noun
within the \synhead{DP}, and the \ird{-ě} form modifies the event as a manner
adverb. Thus \ird{skorn-} means \trsl{be quick}, \ird{skorná} means \trsl{quick}
as a noun modifier, and \ird{skorně} means \trsl{quickly} as a clausal adjunct.
The syntactic position of these derived manner adverbs in the preverbal field is
discussed in §\,\ref{sec:lex-adv}; compare also the simultaneous / manner
converb in \ird{-ěc} in Chapter~\ref{ch:complex}, which remains a non-finite
verbal form rather than a lexical adverb.


A property of a participant can remain constant while the manner of a
particular event varies. Attributive and adverbial \ird{skorn-} therefore
support distinct descriptions:

\pex\label{ex:deriv-slot-property-manner}
\a
\begingl
\gla skorná \textnormal{[[runner]]}\textsubscript{\Dir{}} \textnormal{[[arrive]]}\textsubscript{\Pf{}}.//
\glb quick.\textsc{attr} runner.\Dir{} arrive-\Pf{}//
\glft \trsl{A fast runner arrived.}
\lingcontext{the runner is known to be fast; this arrival need not have been rapid}//
\endgl
\a
\begingl
\gla \textnormal{[[runner]]}\textsubscript{\Dir{}} skorně \textnormal{[[arrive]]}\textsubscript{\Pf{}}.//
\glb runner.\Dir{} quick.\textsc{adv} arrive-\Pf{}//
\glft \trsl{A runner arrived quickly.}
\lingcontext{the arrival was rapid; the runner need not be characteristically fast}//
\endgl
\xe

In (a), the attributive is internal to the nominal headed by `runner';
in (b), the manner derivative modifies the arrival predicate. The change
in attachment leaves the mover direct and the finite verb perfective.

\subsection{Degree and emphatic attribution}\label{sec:stative-degree}
Scalar degree on the stative class is expressed by a closed set of prefixes
that attach to the stative stem before any derivational material, including
the verbalising \ird{-n-} and attributive \ird{-a}/\ird{-á}. The five prefixes are
mutually exclusive and occupy a single slot:

\begin{table}[ht]
  \footnotesize
  \caption{Degree prefixes on stative verbs.}\label{tab:degpfx}
  \begin{tabularx}{\textwidth}{l l L L}
    \toprule
    {\sc prefix} & {\sc value} & {\sc finite} & {\sc attributive} \\
    \midrule
    \ird{po-}  & slight; somewhat      & \ird{pomužna}  & \ird{pomúžna}  \\
    \ird{do-}  & moderate; quite       & \ird{domužna}  & \ird{domúžna}  \\
    \ird{pre-} & strong; very          & \ird{premužna} & \ird{premúžna} \\
    \ird{od-}  & elevated; exceedingly & \ird{odmužna}  & \ird{odmúžna}  \\
    \ird{nad-} & excessive; too        & \ird{nadmužna} & \ird{nadmúžna} \\
    \bottomrule
  \end{tabularx}
\end{table}

\noindent The prefixes are fully productive across the stative class. Because
they precede the stem complex, they appear before \ird{-n-} in finite forms
and before the root in attributives. Each is integrated into the lexical
prosodic word and bears its initial stress: \ird{premužna} and
\ird{premúžna} are stressed on \ird{pre-}, irrespective of the long root
vowel in the root-strong attributive. With external negative \ird{za-}, the stressed
syllable of \ird{zapremužna} remains \ird{pre-}
(§\,\ref{sec:stress}).

Emphatic attribution instead uses the suffix \ird{-ěsk-}, which selects the front grade, so a stative
\ird{-n-} supplies /ɲ/, written \ird{ně} with the following /ɛ/: \ird{mužněská}.
It is distinct from the degree prefix system and does not occupy the same
grammatical slot. It is restricted to the suffix-strong \ird{-á} attributive.
Class~III membership does not license this separate emphatic construction:
neither \ird{všihn-} nor \ird{žbirn-} receives an emphatic principal
part merely by joining the soft class. Its listed quantity is preserved in
\ird{mužněská} and \ird{premužněská}; attributive selection does not reach through
\ird{-ěsk-} to lengthen the root. The construction carries a marked, heightened, or literary flavour rather than a simple
scalar value --- closer in function to an emphatic particle than to a degree
adverb. The contrast with a prefixed form is one of register and framing
rather than degree alone: \ird{premúžna} states a degree as a plain fact;
\ird{mužněská} colours the attribution, implying that the beauty is
remarkable, conspicuous, or worth dwelling on. In ordinary speech
\ird{-ěsk-} attributives may sound affected; in verse, formal oratory, and
deliberate rhetorical prose they are appropriate and expected.

The suffix is not fully productive. It is most natural with roots that carry
inherent aesthetic or evaluative content:

\pex
\a \ird{mužněská} \trsl{strikingly beautiful; beautiful indeed} \lingcontext{natural in literary register}
\a \ird{hradněská} \trsl{truly good; good in the full sense} \lingcontext{natural in literary register}
\xe

\noindent With dimensional or manner-denoting roots the formation is possible
but tends to sound stilted or unidiomatic outside verse:

\pex
\a \ird{svorněská} \trsl{(strikingly big)} \lingcontext{odd in prose; possible in elevated style}
\a \ird{skorněská} \trsl{(remarkably quick)} \lingcontext{marginal; feels forced outside verse}
\xe

\ird{-ěsk-} may co-occur with a degree prefix. The two are additive rather
than competing: the prefix supplies the scalar value and \ird{-ěsk-} supplies
the rhetorical register.

\pex
\a \ird{premužněská} \trsl{very beautiful} \lingcontext{emphatic; degree stated and dwelt upon}
\a \ird{odmužněská} \trsl{exceedingly beautiful} \lingcontext{elevated register; characteristic of panegyric}
\a \ird{nadmužněská} \trsl{excessively, ruinously beautiful} \lingcontext{ironic or elegiac; beauty to a fault}
\xe

\noindent Such combinations are licensed but marked; they are characteristic
of formal verse and deliberate oratorical prose. In ordinary speech, a
prefixed form without \ird{-ěsk-} is the normal choice.

\subsection{Nominalisation and aspectual attribution of statives}

The participant nominaliser \ird{-ou} selects the reduced aspectual
predicate and abstracts over its independently selected participant
(§\,\ref{sec:nominalised}). With a simple intransitive state this is the
holder: continuous \ird{mužnova} means \trsl{a beautiful one}, not
\trsl{beauty}. With an inchoative, \ird{svornélikou} denotes the one
who underwent becoming bigger or big under the supplied scalar standard.
It does not denote the change event. Causative or applicative formation can
alter participant selection under the ordinary rules; a state holder is not
a freely selectable additional gap above those constructions.

\pex
\a \ird{Nužo mužna.} \trsl{The city is beautiful.} \lingcontext{state predication}
\a \ird{Nužo mužnova.} \trsl{The city is the beautiful one.} \lingcontext{participant identification}
\a \ird{Bylo svornélek.} \trsl{The child became bigger.} \lingcontext{bounded inchoative; comparative standard}
\a \ird{Bylo svornélikou.} \trsl{The child is the one who became bigger.} \lingcontext{participant nominal}
\xe

A corresponding \ird{-ový} relative contains the same reduced AspP.
Thus \ird{svornélikový byl} is \trsl{a child who became bigger}, and
\ird{mužnávový nuž} is \trsl{a city which had been beautiful}.
A continuous stative relative is licensed specifically under an episodic,
stage-level, temporally delimited or contrastive interpretation:

\pex
\a \ird{múžna nuž} \trsl{a beautiful city} \lingcontext{unmarked property attributive}
\a \ird{mužnový nuž} \trsl{a city which is beautiful on the relevant occasion} \lingcontext{marked stage-level relative}
\a \ird{*mužnový nuž} \lingcontext{intended as an unrestricted synonym of the lexical attributive}
\xe

The simple attributive is a reduced property projection without viewpoint;
the aspectual relative contains StateP under AspP and abstracts over the
selected holder. This larger structure remains smaller than a finite clause.
Indicative Neg, phasal particles, Mood, clausal evidentials and Comm-field
operators remain excluded \emph{inside} the reduced relative and participant
nominal. A matrix phasal operator may scope over a predication containing
the complete DP; that external relation does not introduce a phasal head
under \ird{-ový} (§\,\ref{sec:nominalised}). When clausal evidence is overtly expressed for a psychological state, its
structure must independently license that evidence; Evid cannot be inserted
under a reduced relative.

Individuals, states, changes and results remain distinct denotations. A
participant nominal of a state denotes its selected bearer, a lexical property
nominal denotes the property, and the separately licensed event-nominal
construction denotes an eventuality. Neither event-derived result meaning
nor state abstraction follows automatically from attaching participant
\ird{-ou}. No new productive state/result nominaliser is specified here.
The root-strong, suffix-strong and soft ordinary attributive classes remain
lexical. The established dimensional attributive \ird{četná} selects
long \ird{-á}, while listed soft \ird{všihňa} selects front-grade
\ird{ň} and short \ird{-a}. Unlisted ordinary attributive class
assignments for \ird{ptakn-} and \ird{hvodn-} cannot be inferred
from their aspectual forms.

\section{Verbal nominalisation}\label{sec:nominalised}\label{sec:nm-derived}


The suffix \ird{-ou} occurs in lexical nouns, event nominals, and
aspect-bearing participant descriptions. Lexical \ird{pjaštou} \trsl{food}
denotes a substance conventionally associated with eating, whereas eventive
\ird{pjaštou} \trsl{eating} denotes an eventuality and admits event-related
dependents, as in \ird{Marcí travěte pjaštou} \trsl{Marek's eating of the
bread}. Participant \ird{pjáščikou} \trsl{the one who ate} denotes an
individual characterised by participation in an event. The shared exponent
therefore does not entail a single internal structure. Event, result, and
participant readings commonly coexist in nominalisation systems
\cite{lieber2018}; the structures distinguished below constitute the analysis
of Iridian adopted here.

Lexical nominalisation combines a root or a small verbal constituent with
$n_{\mathrm{lex}}$, whose interpretation may be stored with the derivative.
Event nominalisation combines an event predicate with $n_{\mathrm{evt}}$,
without selecting an individual participant as its referent. Its established
forms do not require the five-way indicative aspect distinction. Participant
$n_{\mathrm{part}}$ selects \synhead{AspP}, containing
the inner valency domain, optional \synhead{CausP}, and the outer
orientation head.
Embedded Asp retains viewpoint and participant-selection features, although
finite case valuation is inactive: the nominal domain licenses the indirect
case on its dependents (§\,\ref{sec:nm-derived-gerund}). Lexical agent nouns
in \ird{-(h)ár} and \ird{-ist} (§\,\ref{sec:lexical-agent-nouns}) classify a
person by a characteristic activity or occupation; aspect-bearing \ird{-ou}
forms describe a participant in an aspectually construed situation. This
separation accords with Baker and Vinokurova's distinction between lexical
agent nominalisation and the richer verbal structure of reduced relatives
\cite{bakervinokurova2009}.

\subsection{Event nominals and their dependents}
\label{sec:nm-derived-event}
\label{sec:nm-derived-gerund}

An event nominal may occur without overt participants or may be expanded by
indirect-case dependents specifying the actor and theme. These are unexpanded
and expanded uses of the event construction; omission of a dependent does not
turn it into a lexical result noun. The terms do not presuppose that the number
of pronounced dependents diagnoses the theoretical distinction between simple
and complex event nominals.

\pex
\a \ird{pjaštou} \trsl{food} \lingcontext{lexical substance noun}
\a \ird{pjaštou} \trsl{eating} \lingcontext{unexpanded event nominal}
\a \ird{Marcí travěte pjaštou} \trsl{Marek's eating of the bread}
   \lingcontext{expanded event nominal}
\a \ird{Jancí podohletou} \trsl{Janek's remembering / the remembering of Janek}
\xe

The indirect case neutralises the distinction between actor and theme within
the nominal domain. Actor-before-theme order favours the corresponding
interpretation of \ird{Marcí travěte pjaštou}; anchoring the theme may support
that reading when its discourse referent is already recoverable. Neither the
case ending nor \ird{-to} uniquely assigns a thematic role, and a single
dependent such as \ird{Jancí} may accordingly have more than one
interpretation. A possessor shares the indirect case but need not be an actor
or theme, while a postpositional adjunct supplies its own relation. Nominal
case neutralisation has a close comparison in Basque derived event nominals,
whose external and internal arguments may both be genitive, although larger
nominalised clauses retain a different case pattern \cite{artiagoitia2023}.
The comparison supports treating nominal case as a separate domain; it does
not determine which verbal layers Iridian event nominals contain.

\subsection{Event nominals as temporal complements}
\label{sec:nm-derived-temporal}

Event \ird{pjaštou} takes indirect \ird{pjaštově} when the complete
nominal phrase is selected by a temporal postposition
(§\,\ref{sec:nm-postp-temporal}). In constructed
\ird{Marcí travěte pjaštově z} \trsl{since/after Marek's eating of the
bread}, the event head bears the case required by \ird{z}, while
\ird{Marcí} and \ird{travěte} retain their nominal indirect forms. The
actor-before-theme interpretation is supported by the context and ordering;
case itself does not remove the role ambiguity described above. The temporal
construction introduces no new subject case or obligatory pivot identity
between the nominal event and the matrix clause.

The specified singular paradigm also supplies anchored indirect
\ird{pjaštověte}. Anchoring can identify a discourse-relevant event without
assigning it finite tense or making it actual in every interpretive context.
Temporal reference comes from the containing predication and discourse,
including an established episodic occasion or an identifiable projected event.
The temporal postposition then selects the boundary or interval described in
§\,\ref{sec:temporal-event-nominals}; completion and onset are not additional
aspectual exponents on the event nominal.

Neither temporal selection nor expanded argument expression enlarges the
nominal into a finite clause. Independent finite viewpoint, mood, verbal
negation, and clausal evidential suffixes are not licensed inside this event
nominal; matrix operators retain their external scope. Event-denoting
\ird{pjaštově} consequently supplies the temporal complement without
borrowing the finite morphology or participant orientation of a distinct
aspect-bearing participant nominal.

\subsection{Aspect and participant selection}\label{sec:nm-derived-participant}
Participant nominalisation targets the participant privileged by the complete
voice--aspect configuration, rather than assigning an invariant orientation to
a lexical transitive stem. With a transitive base in the perfective or
retrospective, antipassive Voice privileges the agent, while \ird{-n-}
privileges the patient. In the continuous, progressive, and contemplative,
active Voice already privileges the agent, and patient-oriented \ird{-n-}
yields the passive configuration (§\,\ref{sec:vm-n}). Antipassive morphology
remains excluded from these three aspects, including under nominalisation.

\begin{table}[htbp]
\footnotesize
\caption{Participant nominals of \ird{pjašt-} \trsl{eat}. English glosses
in the text give customary contextual readings.}\label{tab:participant-nominals}
\begin{tabularx}{\textwidth}{l L L}
\toprule
Aspect & Agent-oriented & Patient-oriented \\
\midrule
Perfective & \ird{pjáščikou} & \ird{pjaštníkou} \\
Retrospective & \ird{pjáščavou} & \ird{pjaštnávou} \\
Continuous & \ird{pjaštova} & \ird{pjaštnova} \\
Progressive & \ird{pjaštímou} & \ird{pjaštnímou} \\
Contemplative & \ird{pjaštážou} & \ird{pjaštnážou} \\
\bottomrule
\end{tabularx}
\end{table}

The perfective selects quantity-preserving bound
\ird{-ik-}/\ird{-ík-} before nominalising \ird{-ou}.
Thus short \ird{pjášček} corresponds to \ird{pjáščikou},
whereas suffix-strong \ird{pjaštník} gives \ird{pjaštníkou}.
Retrospective \ird{-av}/\ird{-áv} similarly gives
\ird{-avou}/\ird{-ávou}. Continuous Asp together with
$n_{\mathrm{part}}$ has joint direct-singular \ird{-ova}.
Progressive \ird{-ime}/\ird{-íme} corresponds to
\ird{-imou}/\ird{-ímou}, preserving the selected quantity.
Contemplative \ird{-áž}/\ird{-až} combines with \ird{-ou}
as \ird{-ážou}/\ird{-ažou}. These Vocabulary Insertion conditions specify constructional
allomorphy; neither \ird{-ova} nor \ird{-imou} follows from a general
vowel-deletion or root-raising rule. Acute accents retain their ordinary value
as quantity marks rather than stress marks (§\,\ref{sec:stress}).

Perfective \ird{pjáščikou} describes the agent of a bounded eating event,
and retrospective \ird{pjáščavou} an agent whose participation precedes the
relevant reference point. Their customary translations \trsl{the one who ate}
and \trsl{the one who has eaten} do not add grammatical tense. Continuous
\ird{pjaštova} permits habitual or ongoing interpretation, progressive
\ird{pjaštímou} profiles participation during the event's development, and
contemplative \ird{pjaštážou} describes a prospective or intended
participant. The latter may denote an existing individual although the eating
has not occurred.

For the perfective agent form, a schematic interpretation is
\[
 \lambda x\,\exists e\,[\operatorname{eat}(e)\land
 \operatorname{Agent}(e,x)\land\operatorname{PF}(e,R)],
\]
where $R$ is the contextual reference frame. Eventive \ird{pjaštou}, by
comparison, supplies a property of eating events rather than a property of
eaters. The existential closure illustrated for the perfective does not extend
to an assertion of an actual event with the contemplative. Neither
$n_{\mathrm{part}}$ nor \ird{-ou} encodes uniqueness or definiteness;
anchoring and the construction in which the nominal occurs determine its
referential interpretation.

Applicative nominalisation selects the applied participant, as in
\ird{pjaštébikou} \trsl{the one for whose benefit eating took place}.
Causativisation establishes a higher orientation domain, so that participant
nominalisation of an active causative selects the causer in NOM aspects and
the retained theme, otherwise the applied participant, otherwise the causee in
ERG aspects. A NOM-domain causative bearing passive \ird{-an-} selects the
same embedded participant hierarchy instead of the causer
(§\,\ref{sec:vm-n}). The suffixal sequence is root--inner valency--causative--orientation--Asp--nominalisation;
the nominaliser adds no further promotion operation. Thus an inner applicative
need not supply the referent of a causative participant nominal.

\pex
\a \ird{pjáščikou} \trsl{the one who ate}
\a \ird{pjaštova} \trsl{the one who eats}
\a \ird{pjaštníkou} \trsl{that which was eaten}
\a \ird{pjaštébikou} \trsl{the one for whose benefit someone ate}
\xe

The nominal relational head licensing an indirect theme in
\ird{lúbcí pjáščikou} adds a restriction to the event associated with the
participant. It does not restore an overt patient within antipassive
\synhead{VP}. The same participant description occurs in an identificational
predicate, as in \ird{Janek lúbcí pjáščikou}, and in an existential pivot,
as in \ird{Lúbcí pjáščikou jeca}; the difference concerns external
predication rather than nominal derivation (§\,\ref{sec:exst-vs-ident}).
The upper selection boundary excludes internal clausal negation
(§\,\ref{sec:nominalised}). Matrix \ird{Lúbcí pjáščikou zada}
\trsl{Nobody ate the apple} denies satisfiers of a positive description;
\ird{Janek lúbcí pjáščikou staly} rejects Janek's identification with its
participant. Neither means \trsl{someone who did not eat the apple}.

\subsection{The nominal case domain}
A non-privileged participant may be expressed by an indirect-case nominal
dependent, whose relation to the embedded event is licensed in the nominal
domain. In the analysis adopted here, a silent nominal relational head
$\mathrm{R}_{\theta}$ introduces the dependent and adds an event relation
such as Agent or Theme to the participant description before event closure.
Its projection adjoins to the event predicate within the nominaliser's
complement, preserving the selected \synhead{AspP} category in participant
nominals. The relational head licenses an indirect \synhead{KP}; the
dependent's K head realises case, as in other nominal dependency relations
(§\,\ref{sec:nm-poss-struct}). This relation neither assigns a finite verbal
case nor projects another core argument inside \synhead{VP}. Thus a patient
dependent of an antipassive-based nominal
restricts the eating event by Theme$(e,y)$, although antipassive
\synhead{VoiceP} itself contains no overt patient argument. This constructional
relation draws on the lexical event meaning; it does not reverse
antipassivisation. Finite \ird{*Marek travě pjášček} consequently remains
ungrammatical (§\,\ref{sec:vm-antipassive-desc}). The restriction is specific
to Iridian finite antipassives: cross-linguistically, the availability of an
oblique patient varies among antipassive constructions \cite{heaton2020}.

A thematic dependent differs from an ordinary possessor, which relates an
individual to the nominal referent without necessarily identifying an event
role, and from an adjunct such as an instrumental postpositional phrase,
whose relation is contributed by the postposition. All may contain indirect
case, so that case alone identifies neither their interpretation nor their
position in the nominal structure.

\subsection{Inflection, number, and anchoring}
\label{sec:nm-derived-inflection}

The singular \ird{-ou}/\ird{-ov-} alternation is specified for perfective
participant and event nominals. Table~\ref{tab:ou-inflection} gives the three specified singular cells,
with anchored forms after a slash. No plural cell is specified for this
deverbal class; this gap establishes neither zero exponence nor ungrammaticality.

\begin{table}[htbp]
\footnotesize
\caption{Specified singular inflection of event and participant nominals;
plain / anchored forms.}\label{tab:ou-inflection}
% Continuous forms are summarised in the note below the singular paradigm.
\label{tab:ova-inflection}
\begin{tabularx}{\textwidth}{l L L}
\toprule
Case--number & Participant \trsl{one who ate} & Event \trsl{eating} \\
\midrule
Direct singular & \ird{pjáščikou / pjáščikovat} & \ird{pjaštou / pjaštovat} \\
Transitive singular & \ird{pjáščikovy / pjáščikovyte} & \ird{pjaštovy / pjaštovyte} \\
Indirect singular & \ird{pjáščikově / pjáščikověte} & \ird{pjaštově / pjaštověte} \\
\bottomrule
\end{tabularx}
\par\smallskip
\begin{minipage}{\textwidth}
Continuous participant: only unanchored direct singular \ird{pjaštova}
\trsl{one who eats} is specified; all other cells remain unspecified.
\end{minipage}
\end{table}

\label{sec:nm-def-contraction}
The direct anchored \ird{-ovat} form is a construction-conditioned
exponent of the nominal stem, case, and anchoring features. It is not the
output of a general phonological contraction of \ird{-ou} with \ird{-to}.
Class~V transitive \ird{ránou} therefore retains ordinary \ird{ránouto}.

Anchoring a dependent such as \ird{travěte} in
\ird{Marcí travěte pjaštou} marks that dependent as discourse-anchored.
It may help establish the intended interpretation when the context identifies
the bread as the theme, but it does not assign thematic role. The \ird{-ovat}
allomorph of the containing deverbal noun is a separate morphological
operation (§\,\ref{sec:nm-derived-gerund}).

Only the unanchored direct singular \ird{pjaštova} \trsl{one who eats}
is specified for the continuous participant construction. Its ending \ird{-ova} is a joint exponent of Asp and nominal
features, rather than evidence that this construction belongs to ordinary
Class~III \ird{-a} stems. Its obliques and anchored direct form are
unspecified, as are plural exponents for the deverbal class generally. The
\ird{-ou}/\ird{-ov-} singular scheme extends to the other
\ird{-ou}-final aspectual participant forms, without establishing a complete
number paradigm. Consequently, nominal external distribution alone does not
license unrestricted claims about their plural formation.

An event plural, if available, would require countable event units, whereas
lexical \trsl{food} has the mass interpretation described for substance nouns
(§\,\ref{sec:nm-derived-event}). The morphology presently specified does not
determine a distinction among event tokens, repeated occurrences, and event
kinds. Anchoring identifies a discourse-relevant event, participant, or
substance according to the noun's interpretation; it neither changes lexical
\trsl{food} into an argument-taking event nominal nor imposes uniqueness on
every participant description.

\subsection{Negation and the upper boundary}
The aspect-bearing participant nominal and the corresponding \ird{-ový}
relative select actual \synhead{AspP}: viewpoint/Asp.Align together with
orientation \synhead{VoiceP}, optional \synhead{CausP}, and inner
Voice/Appl and lexical structure. The selection boundary is below indicative
\synhead{NegP}, \synhead{PhasP}, the Comm-field, \synhead{MoodP}, and both
clausal evidential layers. None of those higher operators belongs inside the
reduced participant description. Internal verbal \ird{za-} is therefore
excluded; it cannot acquire scope inside the description merely by attaching
to its surface stem. A negative non-indicative suffix likewise cannot enlarge
the nominaliser's complement into a Mood domain.

The structural contrast is
$[n_{\mathrm{part}}\ [\mathrm{AspP}\ \ldots]]$ or
$[\mathrm{RelP}\ \text{\ird{-ový}}\ [\mathrm{AspP}\ \ldots]]$, versus excluded
$*[n_{\mathrm{part}}/\mathrm{Rel}\ [\mathrm{PhasP}/\mathrm{Comm}/
\mathrm{MoodP}/\mathrm{EvidP}/\mathrm{NegP}\ \ldots]]$.
Thus \ird{lúbcí pjáščikou} describes an eater of the apple, and
\ird{lúbcí pjáščikový pres} a man satisfying that same positive description;
neither contains an internal reading \trsl{still/already/only/not eating the
apple}. A matrix phasal or comment operator may instead take scope over a
predication containing the complete nominal or a DP containing the relative.
That external scope does not insert the operator under the reduced head and
does not by itself license a new DP-internal modifier construction. Negation with
\ird{staly} rejects the matrix predication, while \ird{zad-} denies the
existence of an individual satisfying the positive participant description.
Neither expresses \trsl{the one who did not eat}. A finite negative clause
provides the established way to assert non-participation; no internally
negative \ird{-ou} or \ird{-ový} paradigm is licensed here. This restriction
concerns these reduced constructions: it does not prohibit negation in a
finite relative, whose independently licensed finite structure can contain
clausal negation. It supplies no new finite-relative formation rule. Nor is
exclusion of negation a universal property of nominalisation; the smaller
lexical and event nominalisations described here simply supply no productive
clausal negation domain. Strict negative concord remains a property of the
appropriate clausal domain, rather than a mechanism for extending participant
nominalisation above \synhead{AspP}.

\subsection{A possible development of the nominalising family}
\label{sec:nm-derived-history}

A reconstructed deverbal \ird{*-ov-} formation could provide a common source
for event and lexical nouns, with later extension to participant descriptions
in voice-marked structures and thence to adnominal relatives. Positional /v/
vocalisation offers a possible phonetic basis for direct \ird{-ou} beside
prevocalic \ird{-ov-}: modern /v/ has a coda [w] realisation
(Chapter~\ref{ch:phon}). Modern spelling nevertheless preserves \ird{v}
in such contexts, so this reconstruction would also require conventionalisation
of the specialised \ird{-ou} spelling. Anchored \ird{-ovat} would then
involve morphological remodelling on \ird{-ov-}, rather than automatic
synchronic contraction. This sequence is a possible history of the family,
not established Iridian historical evidence; the contemporary structures and
allomorphs do not depend on it. A comparative model for the connection between
role-oriented nominalisation and relatives is the nominalisation-based
Austronesian pattern discussed by Kaufman \cite[§3.4]{kaufman2026}; that
comparison does not establish the proposed Iridian chronology.

\FloatBarrier

\section{Compounding}\label{sec:compound-prosody}

A compound integrated as one lexical prosodic word receives primary stress on
its first syllable, after the stem grades and syllable structure of the complete
form have been established (§\,\ref{sec:stress}). Subsidiary prominence
falls on the third and successive odd-numbered nonfinal syllables of that
complete domain, independently of the boundary between roots
(§\,\ref{sec:secondary-stress}); the second root does not automatically
receive a separate stress. Where the components remain separate prosodic
words, each retains its own initial primary stress and internal footing,
while phrasal prominence may favour one component
(§\,\ref{sec:phrasal-prominence}). The presence of two roots alone
establishes no recursive stress domain or persistent constituent accent.
